\subsection{IDEALS}

It follows from identity (3) that in a Lie algebra \(\mathfrak{g}\) there is no distinction between left ideals and right ideals, every ideal being two-sided. We therefore speak simply of ideals.

*Example. Let G be a Lie group, g its Lie algebra and H a Lie subgroup of G. Every left invariant vector field on H defines canonically a left invariant vector field on G, whence there is a canonical injection of the Lie algebra h of H into g; h is identified with a Lie subalgebra of g under this injection. If H is normal in G, the canonical image of h in g is an ideal of g.*

An ideal of g is a submodule of g which is stable under the inner derivations of g.

DEFINITION 4. A submodule of g which is stable under every derivation of g is called a characteristic ideal of g.

PROPOSITION 2. Let g be a Lie algebra, a an ideal (resp. a characteristic ideal) of g and b a characteristic ideal of a. Then b is an ideal (resp. a characteristic ideal) of g.

Every inner derivation (resp. every derivation) of g leaves a stable and induces on a a derivation and hence leaves b stable.

Let \(\mathfrak{g}\) be a Lie algebra. If \(\mathfrak{a}\) and \(\mathfrak{b}\) are ideals of \(\mathfrak{g}\) , \(\mathfrak{a} + \mathfrak{b}\) and \(\mathfrak{a} \cap \mathfrak{b}\) are ideals of \(\mathfrak{g}\) .

Let \(\mathfrak{a}\) and \(\mathfrak{b}\) be two submodules of \(\mathfrak{g}\) . By an abuse of notation, the submodule of \(\mathfrak{g}\) generated by the elements of the form \([x,y](x\in\mathfrak{a},y\in\mathfrak{b})\) is denoted by \([\mathfrak{a},\mathfrak{b}]\) . We have \([\mathfrak{a},\mathfrak{b}] = [\mathfrak{b},\mathfrak{a}]\) by identity (3). If \(z\in\mathfrak{g}\) , \([z,\mathfrak{a}]\) , or \([\mathfrak{a},z]\) , denotes the submodule \([\mathrm{K}z,\mathfrak{a}] = (\mathrm{ad}\,z)(\mathfrak{a})\) .

PROPOSITION 3. If a and b are ideals (resp. characteristic ideals) of g, {[}a, b{]} is an ideal (resp. a characteristic ideal) of g.

Let D be an inner derivation (resp. a derivation) of g. If \(x \in a\) and \(y \in b\) , then

\[
\mathrm{D} ([ x, y ]) = [ \mathrm{D} x, y ] + [ x, \mathrm{D} y ] \in [ \mathfrak {a}, \mathfrak {b} ].
\]

Hence the proposition.

If \(\mathfrak{a}\) is a submodule of \(\mathfrak{g}\) , the set of \(x \in \mathfrak{g}\) such that (ad \(x\) ). \(\mathfrak{a} \subset \mathfrak{a}\) is a subalgebra \(\mathfrak{n}\) of \(\mathfrak{g}\) called the normalizer of \(\mathfrak{a}\) in \(\mathfrak{g}\) . If further \(\mathfrak{a}\) is a subalgebra of \(\mathfrak{g}\) , then \(\mathfrak{a} \subset \mathfrak{n}\) and \(\mathfrak{a}\) is an ideal of \(\mathfrak{n}\) .

